\documentclass[11pt,a4paper]{article}

% Packages
\usepackage[utf8]{inputenc}
\usepackage[margin=1in]{geometry}
\usepackage{graphicx}
\usepackage{amsmath}
\usepackage{hyperref}
\usepackage{float}
\usepackage{subcaption}
\usepackage{listings}
\usepackage{color}
\usepackage{booktabs}


\definecolor{dkgreen}{rgb}{0,0.6,0}
\definecolor{gray}{rgb}{0.5,0.5,0.5}
\definecolor{mauve}{rgb}{0.58,0,0.82}
\lstset{frame=tb,
  language=C,
  aboveskip=3mm,
  belowskip=3mm,
  showstringspaces=false,
  columns=flexible,
  basicstyle={\small\ttfamily},
  numbers=none,
  numberstyle=\tiny\color{gray},
  keywordstyle=\color{blue},
  commentstyle=\color{dkgreen},
  stringstyle=\color{mauve},
  breaklines=true,
  breakatwhitespace=true,
  tabsize=3
}

\title{Lab 2: Energy Efficient Displays}
\author{Your Name \\ Student ID}
\author{Your Name \\ Student ID}

\date{\today}

\begin{document}


\maketitle


\section{Introduction}
% General introduction about the thesis%

\section{Background}
% Go into the detail of the math especially for the convolution
% Also the sparsity (not implemented yet)

\subsection{Spiking Neural Networks}

\subsubsection{Hard vs Soft Reset}

\subsection{Spiking Convolutional Neural Networks}
% Check the 13-17/04 Week for detail
Spiking Convolutional Neural Networks are SNNs with convolutional connections. 
One of the best networks that achieves 99.1\% accuracy on the MNIST dataset is a Spiking ConvNet, which is trained as an ANN and then transformed into a Spiking ConvNet in Diehl et al.~\cite{diehl2015fast}.
This work also has the Spiking fully connected network version, which has a 98.6\% accuracy with the MNIST dataset. 
These results show us that the convolutional connectivity achieves a higher accuracy. 
Thus, making the convolutional support for the C code generation crucial.


The C code implementation is done with the 



\subsubsection{Backward Pass}

\subsubsection{Forward Pass (Inference)}




\section{Related Work}
% Explain the Simone's work

\section{Methodology and Pipeline Overview}
% Will explain what methods are being used like to test golden ref, speed up calculation etc.
% Will include baseline vs contributions (and maybe its table)%

\section{Optimized C Code Generation}

The basis is the nir\_to\_c\_generator.py script, which generates optimized C code for Cortex-M7 ARM MCUs. 
In the previous version, the capabilities include linear and recurrent connected neural networks with a hard-reset (Reset-to-value) mechanism. 
With my work, the soft-reset (reset-by-subtract) mechanism, convolutional (conv2d) support, and sparsity optimization were introduced.

TODO: Rephrase the sparsity part after implement it.


\subsection{Soft Reset}
The first thing to implement is Soft Reset (Reset-by-subtraction), 
The previous work had only the hard reset (Reset-to-value) and lacked the soft reset mechanism.

First of all, the soft-reset mechanism introduced here follows the same approach as in Eshraghian et al. \cite{eshraghian2023snntorch}.
Also, it is implemented in the \texttt{snnTorch} Python library in the same way, enabling the Python golden reference to match this work. 
The mathematical representation of the soft-reset mechanism is given in the last part of Equation~\ref{eq:soft_reset}.

\begin{equation}\label{eq:soft_reset}
U[t] = \underbrace{\beta U[t - 1]}_{\text{decay}} + \underbrace{W X[t]}_{\text{input}} - \underbrace{S_{\text{out}}[t - 1]\theta}_{\text{reset}}
\end{equation}


In Figure~\ref{fig:soft-reset-c}, the soft-reset implementation in C is given.
First, the current timestep membrane update is being calculated with vectorized ARM functions. 
After the membrane update, the spike check is done. 
While comparing the potential to the threshold, if there is a spike, 
the threshold value is saved to be subtracted at the next timestep. 

\begin{figure}[h]
\centering
\begin{lstlisting}

    // Membrane Update
    q15_t temp1[num_neurons], temp2[num_neurons];

    arm_mult_q15(membrane_potentials, decay_factors, temp1, num_neurons);
    arm_add_q15(temp1, weighted_inputs, temp2, num_neurons);
    arm_sub_q15(temp2, reset_values, membrane_potentials, num_neurons);

    // Spike check, then store reset_value for the NEXT step
    for (uint16_t i = 0; i < num_neurons; i++) {
        if (membrane_potentials[i] > thresholds[i]) {
            output_spikes[i] = 1;
            neurons[i].reset_value = thresholds[i];   // subtract next step
        } else {
            output_spikes[i] = 0;
            neurons[i].reset_value = 0;
        }
        neurons[i].membrane_potential = membrane_potentials[i];
    }

\end{lstlisting}
\caption{Soft-Reset Implementation in C}
\label{fig:soft-reset-c}
\end{figure}





\subsection{Convolutional Support}

First of all a basic convolutional support added. Its pseudo-code shown in the Figure~\ref{}.
% Add pseudo-code
The convolution calculation is done with a sliding window in a deductive manner, from the output neuron to the input neuron.
For every output neuron, the sliding window is being calculated first


\subsection{Sparsity-Aware Inference}


\section{Embedded Deployment}

\subsection{CMSIS-DSP and CMSIS-NN}

\subsection{UART DMA}
% Check the 13-17/04 week for the details.

\subsection{Quad-SPI Flash Chips}
% Week 27-01/05 

\subsection{}


\section{Model Training and Testing}
% I can explain the how I trained the models and how I tested them 
% ???????????????????
% Should I explain the nir problem and manual graph generation in here, or where?

\subsection{Braille}
% Same dataseet used by Simone's.
% to check if the soft-reset implementation is working a network without a recurrent connection is created.
% Check 06-10/04 week for more detail

\subsection{ST-MNIST}
% Explain the ST-MNIST
% Cite this paper https://arxiv.org/abs/2005.04319
% Explain the network and its training
% 20-24/04 for more detail

% MAYBE explain the nir graph generation???????


\section{Experimental Results}

\section{Discussion}


\cite{eshraghian2023snntorch}



\appendix
\section{Appendix}


\bibliographystyle{unsrt} 
\addcontentsline{toc}{chapter}{References}
\bibliography{references}




\begin{figure}[h]
\centering
\begin{lstlisting}
 Example Code
\end{lstlisting}
\caption{Transformation function definitons}
\label{fig:transformation_defs}
\end{figure}


\graphicspath{{images}}
\begin{figure}[H]
  \centering
  %\includegraphics[width=0.7\textwidth]{Pareto_of_default_image_provided.png}
  \caption{Pareto graph of different image transformation techniques}
  \label{fig:pareto_graph1}
\end{figure}



\begin{table}[htbp]
\centering
\caption{Power saving statistics for each transformation}
\label{tab:power_saving_stats}
\begin{tabular}{lrrr}
\toprule
\textbf{Transform} &
\textbf{Mean PS (\%)} &
\textbf{Min PS (\%)} &
\textbf{Max PS (\%)} \\
\midrule
Bright scale & 17.18 & 16.85 & 18.46 \\
Histogram    & -5.53 & -135.09 & 30.56 \\
Hungry blue  & 6.43 & 2.78 & 20.24 \\
\bottomrule
\end{tabular}
\end{table}



\end{document}
